\section{Konvoliuciniai neuroniniai tinklai}
\label{sec:konvoliuciniai_neuroniniai_tinklai}

Šiame skyriuje aprašomos darbe naudojamų konvoliucinių neuroninių tinklų (angl.~\textit{Convolutional Neural Network}, CNN) architektūros bei hiperparametrų reikšmės. Vaizdams taikomos $2$D konvoliucijos, laiko eilutėms – $1$D konvoliucijos, tačiau abu tinklai paveldi tą pačią bazinę klasę \texttt{BaseCNN}, todėl skiriasi tik konvoliucinių, telkimo (angl.~\textit{pooling}) ir paketinio normalizavimo (angl.~\textit{batch normalization}) sluoksnių dimensijos bei įėjimo formos.

\subsection{Architektūra}
\label{subsec:architektura}

Abu tinklai sukonstruoti pagal vienodą šabloną, įgyvendintą bazinėje klasėje \texttt{BaseCNN}. Tinklą sudaro dvi dalys:
\begin{itemize}
    \item \textbf{Požymių išskyrimo blokas} (angl.~\textit{feature extractor}) – kelios konvoliucinių blokų pakopos. Kiekvieną pakopą sudaro: konvoliucinis sluoksnis (\texttt{padding}~$=1$), pasirinktinai paketinio normalizavimo sluoksnis, aktyvacijos funkcija, telkimo sluoksnis ir pasirinktinai \textit{Dropout} sluoksnis.
    \item \textbf{Klasifikatorius} – \texttt{Flatten} sluoksnis, paslėptas pilnai sujungtas (angl.~\textit{fully connected}) sluoksnis, aktyvacijos funkcija, \textit{Dropout} sluoksnis ir išėjimo pilnai sujungtas sluoksnis su tiek neuronų, kiek yra klasių.
\end{itemize}

Skirtingi išvedimo metodai ($\_compute\_flattened\_size$) leidžia teisingai apskaičiuoti suplotos požymių žemėlapio matricos dydį po visų telkimo sluoksnių: vaizdams – $C \cdot (W / s^{L})^2$, laiko eilutėms – $C \cdot \lfloor n / s^{L} \rfloor$, čia $C$ – paskutinio sluoksnio filtrų skaičius, $W$ – įėjimo vaizdo kraštinė, $n$ – įėjimo požymių skaičius, $s$ – telkimo dydis, $L$ – konvoliucinių blokų skaičius.

\subsubsection{Vaizdų tinklas (2D CNN)}
\label{subsubsec:image_cnn}

Klasė \texttt{ImageCNN} naudoja \texttt{nn.Conv2d}, \texttt{nn.BatchNorm2d}, \texttt{nn.MaxPool2d} ir \texttt{nn.Dropout2d} sluoksnius. Įėjimas – $3 \times 64 \times 64$ formato RGB tenzorius (vaizdo kraštinė valdoma konstanta \texttt{IMAGE\_SIZE} $= 64$). Bazinė konfigūracija turi tris konvoliucinius blokus su filtrų skaičiais $(32,\,64,\,128)$, branduolio dydžiu $3 \times 3$ ir telkimo dydžiu $2 \times 2$. Po trijų telkimo sluoksnių požymių žemėlapio kraštinė sumažėja iki $64 / 2^{3} = 8$, todėl po suplojimo gaunamas $128 \cdot 8 \cdot 8 = 8192$ dydžio vektorius, kuris paduodamas į $256$ neuronų paslėptąjį sluoksnį, o galiausiai – į $2$ neuronų išėjimo sluoksnį (klasės \textit{chihuahua} ir \textit{muffin}). Sluoksnių parametrai pateikti \ref{tab:image_cnn_architecture} lentelėje.

\begin{table}[H]
    \centering
    \caption{Vaizdų klasifikavimo tinklo (\texttt{ImageCNN}) bazinės architektūros sluoksniai.}
    \begin{tabular}{|c|l|l|}
        \hline
        Nr. & Sluoksnis & Parametrai / išėjimas \\
        \hline
        $0$ & Įėjimas & $3 \times 64 \times 64$ \\ \hline
        $1$ & \texttt{Conv2d} & $3 \to 32$, branduolys $3$, \texttt{padding}~$=1$ \\ \hline
        $2$ & \texttt{ReLU} & – \\ \hline
        $3$ & \texttt{MaxPool2d} & dydis $2$, išėjimas $32 \times 32 \times 32$ \\ \hline
        $4$ & \texttt{Conv2d} & $32 \to 64$, branduolys $3$, \texttt{padding}~$=1$ \\ \hline
        $5$ & \texttt{ReLU} & – \\ \hline
        $6$ & \texttt{MaxPool2d} & dydis $2$, išėjimas $64 \times 16 \times 16$ \\ \hline
        $7$ & \texttt{Conv2d} & $64 \to 128$, branduolys $3$, \texttt{padding}~$=1$ \\ \hline
        $8$ & \texttt{ReLU} & – \\ \hline
        $9$ & \texttt{MaxPool2d} & dydis $2$, išėjimas $128 \times 8 \times 8$ \\ \hline
        $10$ & \texttt{Flatten} & $8192$ \\ \hline
        $11$ & \texttt{Linear} & $8192 \to 256$ \\ \hline
        $12$ & \texttt{ReLU} & – \\ \hline
        $13$ & \texttt{Dropout} & $p = 0{,}0$ (bazinė) \\ \hline
        $14$ & \texttt{Linear} & $256 \to 2$ \\ \hline
    \end{tabular}
    \label{tab:image_cnn_architecture}
\end{table}


\subsubsection{Laiko eilučių tinklas (1D CNN)}
\label{subsubsec:keystroke_cnn}

Klasė \texttt{KeystrokeCNN} naudoja \texttt{nn.Conv1d}, \texttt{nn.BatchNorm1d}, \texttt{nn.MaxPool1d} ir \texttt{nn.Dropout} sluoksnius. Įėjimas – $1 \times 31$ formato sekos tenzorius ($31$ klaviatūros laikinių požymių, vienas kanalas). Bazinė konfigūracija turi du konvoliucinius blokus su filtrų skaičiais $(64,\,128)$, branduolio dydžiu $3$ ir telkimo dydžiu $2$. Po dviejų telkimo sluoksnių sekos ilgis sumažėja iki $\lfloor 31 / 2^{2} \rfloor = 7$, todėl po suplojimo gaunamas $128 \cdot 7 = 896$ dydžio vektorius, paduodamas į $128$ neuronų paslėptąjį sluoksnį, o galiausiai – į $30$ neuronų išėjimo sluoksnį (po vieną neuroną kiekvienam dalyviui). Sluoksnių parametrai pateikti \ref{tab:keystroke_cnn_architecture} lentelėje.

\begin{table}[H]
    \centering
    \caption{Laiko eilučių klasifikavimo tinklo (\texttt{KeystrokeCNN}) bazinės architektūros sluoksniai.}
    \begin{tabular}{|c|l|l|}
        \hline
        Nr. & Sluoksnis & Parametrai / išėjimas \\
        \hline
        $0$ & Įėjimas & $1 \times 31$ \\ \hline
        $1$ & \texttt{Conv1d} & $1 \to 64$, branduolys $3$, \texttt{padding}~$=1$ \\ \hline
        $2$ & \texttt{ReLU} & – \\ \hline
        $3$ & \texttt{MaxPool1d} & dydis $2$, išėjimas $64 \times 15$ \\ \hline
        $4$ & \texttt{Conv1d} & $64 \to 128$, branduolys $3$, \texttt{padding}~$=1$ \\ \hline
        $5$ & \texttt{ReLU} & – \\ \hline
        $6$ & \texttt{MaxPool1d} & dydis $2$, išėjimas $128 \times 7$ \\ \hline
        $7$ & \texttt{Flatten} & $896$ \\ \hline
        $8$ & \texttt{Linear} & $896 \to 128$ \\ \hline
        $9$ & \texttt{ReLU} & – \\ \hline
        $10$ & \texttt{Dropout} & $p = 0{,}0$ (bazinė) \\ \hline
        $11$ & \texttt{Linear} & $128 \to 30$ \\ \hline
    \end{tabular}
    \label{tab:keystroke_cnn_architecture}
\end{table}



\subsection{Hiperparametrai}
\label{subsec:hiperparametrai}

Visiems eksperimentams bendrosios mokymo hiperparametrų reikšmės yra fiksuotos: paketo dydis (angl.~\textit{batch size}) – $64$, epochų skaičius – $20$, mokymosi greitis (angl.~\textit{learning rate}) – $10^{-3}$, paklaidos funkcija – \textit{Cross Entropy}, optimizavimo algoritmas – \textit{Adam}, branduolio dydis – $3$, telkimo dydis – $2$. Šios reikšmės keičiamos tik tuose tyrimuose, kurie skirti būtent tam hiperparametrui įvertinti. Bendrieji hiperparametrai surinkti \ref{tab:hiperparametrai_bendri} lentelėje, o tyrimuose varijuojamos reikšmės – \ref{tab:hiperparametrai_tyrimai} lentelėje.

\begin{table}[H]
    \centering
    \caption{Bendrosios (fiksuotos) hiperparametrų reikšmės abiem tinklams.}
    \begin{tabular}{|l|l|}
        \hline
        Hiperparametras & Reikšmė \\
        \hline
        Paketo dydis (\textit{batch size}) & $64$ \\ \hline
        Epochų skaičius & $20$ \\ \hline
        Mokymosi greitis (\textit{learning rate}) & $10^{-3}$ \\ \hline
        Optimizavimo algoritmas & \textit{Adam} \\ \hline
        Paklaidos funkcija & \textit{Cross Entropy} \\ \hline
        Aktyvacijos funkcija & \textit{ReLU} \\ \hline
        Branduolio dydis (\textit{kernel}) & $3$ \\ \hline
        Telkimo dydis (\textit{pool}) & $2$ \\ \hline
        \texttt{padding} konvoliucijose & $1$ \\ \hline
        \textit{Dropout} tikimybė (bazinė) & $0{,}0$ \\ \hline
        Paketinis normalizavimas (bazinė) & išjungta \\ \hline
        Pradinė \textit{seed} reikšmė & $42$ \\ \hline
    \end{tabular}
    \label{tab:hiperparametrai_bendri}
\end{table}

\begin{table}[H]
    \centering
    \caption{Tyrimuose varijuojami hiperparametrai ir jų reikšmės.}
    \begin{tabular}{|l|l|l|}
        \hline
        Tyrimas & Vaizdų variantai & Laiko eilučių variantai \\
        \hline
        Architektūra (filtrai) & $(16,\,32)$; $(32,\,64,\,128)$; $(64,\,128,\,256)$ & $(32)$; $(64,\,128)$; $(64,\,128,\,256)$ \\ \hline
        \textit{Dropout} tikimybė & \multicolumn{2}{c|}{$0{,}0$; $0{,}25$; $0{,}5$; $0{,}75$} \\ \hline
        Paketinis normalizavimas & \multicolumn{2}{c|}{išjungta; įjungta} \\ \hline
        Aktyvacijos funkcija & \multicolumn{2}{c|}{\textit{ReLU}, \textit{LeakyReLU}, \textit{Tanh}, \textit{ELU}} \\ \hline
        Optimizavimo algoritmas & \multicolumn{2}{c|}{\textit{Adam}, \textit{SGD}, \textit{RMSprop}, \textit{AdamW}} \\ \hline
    \end{tabular}
    \label{tab:hiperparametrai_tyrimai}
\end{table}


Atsitiktinių skaičių generatoriaus pradinė reikšmė (angl.~\textit{seed}) buvo fiksuota ($seed = 42$) duomenų dalijime ir poaibių formavime, kad eksperimentų rezultatai būtų atkartojami.

